% Insert as A.61.  A.62 constructs the ordinary locus used here.
\begin{statement}{A.61}{Proved}{An ordinary valuation locus in the actual full class-four image}
Put $E_4=\pi_1(U)/\Gamma_5$, with the lower central
series taken in the geometric group.  Let $\mathcal V_4$
be the ordinary completed unit-equation pairs with
valuation patterns $(0,1)$ at two and $(1,0)$ at three,
and, at every $q>3$, one of
\[
 (0,0),\qquad (E_q,0),\qquad(0,E_q),\qquad(-E_q,-E_q),
 \quad E_q=24e_q(q^2-1)(q^3-1).
 \tag{A.61.1}
\]
Here $e_q$ are nonnegative integers forming a family
which converges to zero in $\widehat\Z$.
The image of the actual section map
\[
       \operatorname{Sec}(E_4/G_\Q)
                 \longrightarrow\operatorname{Sec}(E_{\rm ab}/G_\Q)
 \tag{A.61.2}
\]
contains $\mathcal V_4$.  Every class thus obtained has
the distinguished intermediate class-two section $s_*$.
The construction uses the full integral fourth central
layer, including every coefficient prime and its
continuous cohomology.
The explicit pair of A.62 with $C=24$ and $J=6$ belongs
to this locus. Thus a single ordinary infinite-support pair
has an actual full integral class-four lift; its abelian
support and cusp maps exclude rational-point and cuspidal
images, as in A.59.
\end{statement}
\begin{proof}
\emph{The actual fourth central layer and its coordinate maps.}
Use the generators $x,y$ and commutator convention of A.43,
and set
\[
 C=[x,y],\quad D=[C,x],\quad F=[C,y],\qquad
 P=[D,x],\quad Q=[D,y],\quad R=[F,y].
 \tag{A.61.3}
\]
The Hall--Witt identity gives $[D,y]=[F,x]$ modulo
$\Gamma_5$.  Bilinearity of the associated graded
commutators shows that $P,Q,R$ generate
$\Gamma_4/\Gamma_5$.
We verify their independence integrally, rather than
assuming a faithful full profinite Magnus embedding.
Use only the truncated Magnus homomorphism
\[
 x\longmapsto1+X,\qquad y\longmapsto1+Y
       \quad\hbox{in }1+(X,Y)
       \subset\widehat\Z\langle X,Y\rangle/(X,Y)^5.
 \tag{A.61.4}
\]
On the degree-four words $XXXY,XXYY,XYYY$, the
coefficient columns of $P,Q,R$ are
\[
                   \begin{pmatrix}1&0&0\\0&-1&0\\0&0&1\end{pmatrix}.
 \tag{A.61.5}
\]
Thus the surjection from $\widehat\Z^3$ given by
these generators has trivial kernel.  In particular
\[
        \ker(E_4\to E_3)=\Gamma_4/\Gamma_5
                              =\widehat\Z(4)^3.
 \tag{A.61.6}
\]
The twist follows by applying the actual Ihara action
to the weight-four commutators: the degree-one action
is multiplication by $\chi$, hence the degree-four
action is $\chi^4$.

The same coefficient argument, inductively with A.43,
gives the continuous ordered Hall coordinates
\[
                 y^b x^a C^h D^d F^f P^uQ^vR^w.
 \tag{A.61.7}
\]
Existence follows by collecting commutators successively
in each of the four central quotients; independence follows
from the primitive coefficient maps at each weight.
Setting the last three coordinates to zero is an actual
continuous set section of $E_4\to E_3$, over the fixed
rational tangential splitting.  Its defect gives the
continuous obstruction cocycle with coefficients (A.61.6).
The projection of (A.61.7) to $\Z_\ell$ coordinates will
be used below.  No assertion about a faithful Magnus map
on the untruncated full group is needed.

\emph{Unramified action and the exact tame residue.}
Fix $q\ne\ell$ and work in the pro-$\ell$ geometric
quotients $N_i=\Pi_\ell/\Gamma_{i+1}$.
The actual Galois action on $N_2$ is unramified at $q$:
$\chi_\ell$ is unramified, and in the Ihara formula
$g(y)=f_g^{-1}y^{\chi(g)}f_g$ the element $f_g\in\Gamma_2$
is central on $N_2$.  Thus inertia fixes both generators.
This also proves unramifiedness on $N_3$ and $N_4$ by
the following exact comparison of automorphisms.
For $i=2,3$, the kernel of
$\operatorname{Aut}(N_{i+1})\to\operatorname{Aut}(N_i)$
consists of the maps changing each free generator by
a central element of $\Gamma_{i+1}/\Gamma_{i+2}$.
It is
\[
 \operatorname{Hom}\bigl(N_1,\Gamma_{i+1}/\Gamma_{i+2}\bigr),
 \tag{A.61.8}
\]
a torsion-free $\Z_\ell$ module on which conjugation by
arithmetic Frobenius acts by $q^i$.
Indeed the two factors in this conjugation action are
$q^{i+1}$ on the target and $q^{-1}$ on the source.
Once the action on $N_i$ is unramified, inertia maps
homomorphically into (A.61.8).  Wild inertia and the
prime-to-$\ell$ part of tame inertia have zero image.
For a tame $\ell$-generator $T$ and Frobenius $\Phi$,
the relation $\Phi T\Phi^{-1}=T^q$ gives
$(q^i-q)\delta(T)=0$ in (A.61.8).  Torsion-freeness
makes $\delta(T)=0$.  This proves the stated
unramifiedness through class four without assigning
unproved values to higher Ihara coefficients.

Let $s$ be an actual class-three section.  Its restriction
at $q$ therefore factors, for its geometric cocycle and
action, through the tame quotient just used.
Write
\[
 V=s(T)_{\rm geom}=y^n x^m C^t D^dF^f,\qquad
 U=s(\Phi)_{\rm geom}.
 \tag{A.61.9}
\]
The exact class-three tame relation is
\[
                   U\varphi_q(V)U^{-1}=V^q,
 \tag{A.61.10}
\]
where $\varphi_q$ is the actual unramified Frobenius
automorphism of the geometric group.
Let $\widetilde U,\widetilde V$ be the lifts obtained
by (A.61.7).  Then
\[
 \mathcal R_q=\widetilde U\varphi_q(\widetilde V)
                          \widetilde U^{-1}\widetilde V^{-q}
                      \ \in\Gamma_4/\Gamma_5.
 \tag{A.61.11}
\]
For the normalized global defect cocycle $c$, the strict
cochain residue of A.27, now in weight four, is exactly
\[
 \begin{split}
 \psi_q(c)=q^{-1}\biggl(&c(\Phi,T)+c(\Phi T,\Phi^{-1})
          -c(\Phi,\Phi^{-1})
          -\sum_{j=1}^{q-1}c(T^j,T)\biggr)
                   =q^{-1}\mathcal R_q.
 \end{split}
 \tag{A.61.12}
\]
The last equality follows by multiplying the chosen lifts:
the first three terms compare
$\widetilde\Phi\widetilde T\widetilde\Phi^{-1}$
with the lift of $T^q$, and the sum compares
$\widetilde T^q$ with that same lift.
In (A.61.12) the central coordinates are written
additively.  Changing the lift of $T$ by a central
coordinate $z$ changes $\mathcal R_q$ by
$(q^4-q)z$, so the displayed residue differential is
precisely $q^3-1$.
This identifies the actual cocycle residue, not only
its class in a local cohomology group.

\emph{An integral bound from the actual Hall lifts.}
At a prime $q>3$ in (A.61.1), put $E=E_q$.
Ordinary local Kummer coordinates and the class-two
tame relation, as in A.46 and A.58, give
\[
                 m,n\in E\Z_\ell,\qquad
                              t\in(E/2)\Z_\ell.
 \tag{A.61.13}
\]
In a positive valuation pattern the value of $t$ is
$mv$ or $-nu$ with $u,v\in\Z_\ell$; in the negative
pattern it is $\binom m2-mu$.  These are the actual
cochain identities giving (A.61.13).

Set $V_0=y^n x^mC^t$.
Its Magnus coefficients through degree three, except
the constant coefficient, belong to $(E/6)\Z_\ell$.
To check every denominator in this assertion, use
\[
                \binom ak=\frac ak\binom{a-1}{k-1}
                         \quad(a\in\Z_\ell,\ 1\le k\le4).
 \tag{A.61.14}
\]
For $a=m,n$ and $k\le3$ the only losses divide six.
Through degree three, $C^t$ is linear in $t$ because
$C-1$ starts in degree two; (A.61.13) loses at most
two.  Products preserve the resulting integral ideal.

Modulo class three, $D,F$ are central and Frobenius
acts on them by $q^3$.  Substituting
$V=V_0D^dF^f$ in (A.61.10) therefore gives
\[
 U\varphi_q(V_0)U^{-1}V_0^{-q}
                   =D^{-q(q^2-1)d}F^{-q(q^2-1)f}.
 \tag{A.61.15}
\]
Magnus substitution by the actual Frobenius action,
and conjugation by the integral element $U$, preserve
every integral coefficient ideal.  In fact each generator
is substituted by an integral power series with zero
constant term after subtracting one; all inverse group
elements also have integral truncated series.
The left side of (A.61.15) is consequently one modulo
$(E/6)\Z_\ell$.  On the degree-three words $XXY,XYY$,
the coefficient matrix of $D,F$ is
$\operatorname{diag}(-1,1)$.  Thus (A.61.15) proves
\[
                d,f\in\frac{E}{6(q^2-1)}\Z_\ell.
 \tag{A.61.16}
\]
Here division is first interpreted in $\Q_\ell$;
the displayed ideal is integral for (A.61.1).

Through degree four, the nonconstant coefficients of
$V_0$ belong to $(E/12)\Z_\ell$.
Formula (A.61.14) now loses only divisors of
$\operatorname{lcm}(1,2,3,4)=12$ for the powers
of $x,y$.  The only additional term in $C^t$ is
$\binom t2(C-1)^2$, whose coefficients belong to
$(E/4)\Z_\ell$ by (A.61.13).
The powers $D^d,F^f$ remain linear through degree
four because they start in degree three.
Combining this with (A.61.16) shows that the actual
Hall lift $\widetilde V$ in (A.61.11) is one modulo
the coefficient ideal
\[
               J_q=\frac{E}{12(q^2-1)}\Z_\ell
                         =2e_q(q^3-1)\Z_\ell.
 \tag{A.61.17}
\]
Frobenius substitution, integral conjugation, and the
integer $q$th power preserve this ideal.  Hence
$\mathcal R_q-1$ has all its Magnus coefficients
in $J_q$.  It is central by (A.61.10), and the
unimodular matrix (A.61.5) extracts its three actual
integral central coordinates.  Each coordinate of
$\psi_q(c)/(q^3-1)$ therefore belongs to
$e_q\Z_\ell$.

If both abelian valuations at $q$ are zero, the
class-two tame relation gives $t=0$.  Equation
(A.61.15), with $V_0=1$, then gives $d=f=0$.
In this case $\widetilde V=1$ and (A.61.12) is
identically zero.  This also proves that the
strict residue families have no additional support
at ordinary primes of valuation pair $(0,0)$.
The exact truncated computations and ideal checks are
retained in \texttt{class4\_magnus\_certificate.py}
and its output: both primitive coefficient matrices,
the relation $Q=[F,x]$, central-coordinate extraction,
and the degree-three $E/6$ and degree-four $E/12$ bounds.

\emph{The global unramified kernel and the old places.}
The strict localization comparison of A.27, with
coefficient module $\Z_\ell(4)$, has residue complex
\[
 R_\ell\xrightarrow{\ d_4\ }R_\ell,\qquad
                     (d_4z)_q=(q^3-1)z_q.
 \tag{A.61.18}
\]
It gives an exact sequence at the middle term
\[
 H^2_{\rm cts}(G_{\Q,\{\ell,\infty\}},\Z_\ell(4))
 \longrightarrow H^2_{\rm cts}(G_\Q,\Z_\ell(4))
 \longrightarrow R_\ell/d_4R_\ell.
 \tag{A.61.19}
\]
For odd $\ell$ the left group is zero:
\cite[VI, Theorem 8.2 and Examples 10.3]{Weibel2013}
identify it with the $\ell$-primary part of
$K_6(\Z[1/\ell])$, and that part equals the
$\ell$-primary part of $K_6(\Z)=0$ by localization.
At two, the needed precise conclusion is injectivity of
\[
 H^2_{\rm cts}(G_{\Q,\{2,\infty\}},\Z_2(4))
                \longrightarrow H^2(\R,\Z_2(4)).
 \tag{A.61.20}
\]
Indeed \cite[VI, Theorem 9.11, case $n=6$, and
Examples 10.3]{Weibel2013} identify its kernel with
$K_6(\Z[1/2])\{2\}=K_6(\Z)\{2\}=0$.
The real group on the right is $\Z/2$ and is retained
in this argument.

A pair in $\mathcal V_4$ belongs to $\mathcal V_3$:
the latter factorial factors are $e_q(q^3-1)$,
which still converge to zero.  A.58 therefore gives
an actual class-three lift of its $s_*$.
Use A.60 to choose such a lift matching the actual
ordinary point sections at $(2,7),(3,13),(3,2)$
and at $(\infty,2)$.  Its local fourth-layer
obstruction is zero at these places because those
point sections already lift to the full group.
These are all the nonunit tame denominators at two
and three: $2^3-1=7$ and $3^3-1=26$.
For every other coefficient prime at those two
places the corresponding denominator is a unit,
while the coefficient prime equal to the place is
omitted from the residue family.

It follows that every coordinatewise quotient
$\psi_q(c)/(q^3-1)$ is integral.  At the fresh support
primes its entries belong to $e_q\Z_\ell$ by
(A.61.17), and at all remaining primes they are
zero.  They consequently form elements of the
literal module $R_\ell$; the finitely many entries
at two and three do not affect convergence.
Thus all three full obstruction classes have zero
image in (A.61.19), including its continuous
derived-limit contribution.  For odd $\ell$ they
are zero.  At two their real restrictions are zero
by the chosen real point comparison, and (A.61.20)
makes the unramified classes zero as well.

Finally, take the product over coefficient primes
in the actual central extension (A.61.6).
The continuous cochain complex with these product
coefficients is the product of the individual
complexes.  The three obstruction coordinates
therefore have a global continuous primitive,
which corrects the set section (A.61.7) to the
required arithmetic section of $E_4$.
Its class-two image is $s_*$ because all corrections
were in the higher central kernels.
\end{proof}

